\section{Usage notes}\label{usage}

The portal at \url{https://rubalkhali.science/} provides site pages and
a SPARQL editor. The endpoint
\url{https://rubalkhali.science/sparql} accepts SPARQL~1.1 queries.
For a local, low-memory example, load the base, site and XRF modules
listed in Table~\ref{tab:modules} and execute
\path{sparql/field_xrf_site10.rq}. The archived Site~10 result contains
46 rows; compare both the returned process/analyte identifiers and
their values with \path{evidence/competency-query/}. Taxonomic
queries additionally require the taxonomy ABox and its mapping
dependencies. The named-run example retrieves the complete genus
profile for \texttt{ERR16061083}. For all-run exports, execute
\path{scripts/release/kg/taxonomy_all_runs.rq} on a locally loaded
snapshot and allow for a long-running join. The release's \path{REPRODUCE.md} describes payload
retrieval, checksums, environments and workflow commands.

\subsection{Choosing observations for reuse}

Use specimen, collection-visit and assay identifiers when joining
observations. The primary repeated-campaign frame comprises sites
1--60; the source identifiers 61--64 refer to four named catalogue
locations sampled in Trip~1. Several feature profiles share a source
specimen identifier. Analyses should group these profiles by their
documented specimen and processing relationships.

The control registry distinguishes assay, processing stage,
extraction group and identity evidence. The Trip~5 candidate screen
uses 17 extraction blanks linked to 217 biological profiles and
identifies 351 ASVs. The canonical feature table retains these ASVs;
the companion sensitivity analysis applies their removal to the linked
profiles. Trip~4 has six separately denoised extraction blanks and
three batches with unsequenced blanks, leaving complete or partial
coverage gaps at 28 sites. Concentration measurements exist in the
source records; any concentration-based screen requires a matched
specimen--extraction--measurement ledger. Four Trips~2--3 standards
support genus-recovery evaluation. Trips~4--5 have extraction controls
but no sequenced 16S positive-standard profiles.

Archived-soil pH describes specimens measured in CaCl$_2$ under the
admission rules of Section~\ref{methods:ph}. Coverage varies by
campaign and compartment. All Trip~4 assays used archived field
replicate~2. Their supported comparison with community profiles is at
campaign, site and compartment level. The versioned reconciliation
links measurements to the confirmed specimens as described in Methods.

Field and laboratory XRF observations have separate acquisition
provenance and undocumented concentration units. The released values
support within-workflow descriptive comparisons; cross-workflow
calibration requires documented units, paired aliquots and reference
materials. Site~4 field readings include different substrates, so a
site-level summary combines that local heterogeneity with instrument
repeatability.

The PMA table contains three site-compartment groups at two campsites.
Each group has one parent collection tube, from which three separate
slurries supplied treated/untreated pairs. Descriptive paired summaries
retain this preparation structure and compare dye treatment under
matched dark and light exposure.

\subsection{Versioning and supported operations}

The deposited source and checksum manifests identify the inputs used
to generate each module. Preserve these files when reproducing a
reported result. Rebuilding against a newer external ontology requires
identifier and mapping validation, followed by regeneration of the
affected modules and query results. The repository changelog records
source corrections and superseded artifacts.

KG v3.0.0 serves the asserted graph at
\url{https://rubalkhali.science/graph/asserted/3.0.0}.
This graph is the default dataset for portal and SPARQL queries.
Explicit \texttt{FROM}, \texttt{FROM NAMED} and SPARQL protocol
dataset parameters select the dataset for an individual request.
The service evaluates stored assertions directly and rejects
unavailable query modes.

The bundled finite-rule implementation supports the separate
reasoning workflow described in Section~\ref{subsec:owl_rdf}.
Publication of its output requires completion and validation of that
workflow. Changes to released artifacts require a new semantic version.

For reproduction, retain the source package, module manifest,
loader provenance and checksums together. The release instructions
specify the pinned generation environment and coordinate-preserving
import adapter. The asserted N-Quads provides the released union;
the source package regenerates all 16 original module hashes.
Local module loading provides access to fixed inputs independently
of the portal's deployed version. Analyses requiring additional OWL
consequences need a reasoner and validation workload suited to those
constructs.

\section{Data availability}\label{data_availability}

The project uses ENA umbrella study \texttt{PRJEB104209}; its submitted
amplicon run records are associated with \texttt{PRJEB106069}. On
10~September 2026, the anonymous ENA read-run file report for
\texttt{PRJEB106069} returned no public run/file rows. Release of the
raw reads or a complete anonymous reviewer download remains required.
The source and derived-data package is maintained at
\url{https://github.com/bio-ontology-research-group/empty-quarter-ecology-data}.
Its manifests distinguish included files, separately retrieved bulk
payloads and unresolved release gates. KG v3.0.0 and its source package
are available through the versioned website downloads. The broader
dataset's archival accession, sequence access and project-data licence
remain separate publication requirements.

\section{Code availability}\label{code}

The web application and ontology-development code are available at
\url{https://github.com/bio-ontology-research-group/empty-quarter}.
The data repository
\url{https://github.com/bio-ontology-research-group/empty-quarter-ecology-data}
contains the manuscript sources, RDF generators, metadata and
\texttt{main.nf} workflow. RDF generation uses Groovy scripts with
OWL API v5.1.20 \cite{owlapi}; the ELK validation wrapper uses
OWL API v4.5.26. The repository supplies execution profiles, workflow
tests, a pinned Conda recipe and a hash-locked Python environment.
